
\documentclass{article}
\usepackage{colortbl}
\usepackage{makecell}
\usepackage{multirow}
\usepackage{supertabular}

\begin{document}

\newcounter{utterance}

\centering \large Interaction Transcript for game `hot\_air\_balloon', experiment `air\_balloon\_survival\_en\_reasoning off\_hard', episode 1 with gpt{-}5{-}mini{-}2025{-}08{-}07{-}t1.0.
\vspace{24pt}

{ \footnotesize  \setcounter{utterance}{1}
\setlength{\tabcolsep}{0pt}
\begin{supertabular}{c@{$\;$}|p{.15\linewidth}@{}p{.15\linewidth}p{.15\linewidth}p{.15\linewidth}p{.15\linewidth}p{.15\linewidth}}
    \# & $\;$A & \multicolumn{4}{c}{Game Master} & $\;\:$B\\
    \hline

    \theutterance \stepcounter{utterance}  
    & & \multicolumn{4}{p{0.6\linewidth}}{
        \cellcolor[rgb]{0.9,0.9,0.9}{
            \makecell[{{p{\linewidth}}}]{
                \texttt{\tiny{[P1$\langle$GM]}}
                \texttt{Stai partecipando a un gioco di negoziazione collaborativa.} \\
\\ 
\texttt{Insieme a un altro partecipante, dovete accordarvi su un unico set di oggetti da conservare. Ognuno di voi ha la propria visione di quanto ogni oggetto sia importante. Non sai come l'altro partecipante valuta gli oggetti. Inoltre, ti viene fornito lo sforzo che ogni oggetto richiede.} \\
\texttt{Potete accordarvi su un set solo se lo sforzo totale degli oggetti selezionati non supera un limite condiviso:} \\
\\ 
\texttt{LIMITE = 3357} \\
\\ 
\texttt{Ecco i valori di sforzo dei singoli oggetti:} \\
\\ 
\texttt{Sforzo oggetto = \{"A21": 506, "B93": 492, "A11": 338, "B33": 122, "C08": 131, "B39": 909, "B96": 144, "C85": 716, "A82": 263, "B31": 231, "B20": 91, "A60": 651, "B09": 552, "C81": 849, "B21": 720\}} \\
\\ 
\texttt{Ecco la tua visione personale sull'importanza di ogni oggetto:} \\
\\ 
\texttt{Valori di importanza dell'oggetto = \{"A21": 138, "B93": 583, "A11": 868, "B33": 822, "C08": 783, "B39": 65, "B96": 262, "C85": 121, "A82": 508, "B31": 780, "B20": 461, "A60": 484, "B09": 668, "C81": 389, "B21": 808\}} \\
\\ 
\texttt{Obiettivo:} \\
\\ 
\texttt{Il tuo obiettivo è negoziare un set condiviso di oggetti che ti avvantaggi il più possibile (cioè, massimizzi l'importanza totale per TE), rimanendo entro il LIMITEE. Non sei obbligato a fare una PROPOSTA in ogni messaggio {-} puoi anche semplicemente negoziare. Tutte le tattiche sono permesse!} \\
\\ 
\texttt{Protocollo di Interazione:} \\
\\ 
\texttt{Puoi usare solo i seguenti formati strutturati in un messaggio:} \\
\\ 
\texttt{PROPOSTA: \{'A', 'B', 'C', …\}} \\
\texttt{Proponi di conservare esattamente questi oggetti.} \\
\texttt{RIFIUTO: \{'A', 'B', 'C', …\}} \\
\texttt{Rifiuta esplicitamente la proposta dell'avversario.} \\
\texttt{ARGOMENTAZIONE: \{'...'\}} \\
\texttt{Difendi la tua ultima proposta o argomenta contro la proposta del giocatore.} \\
\texttt{ACCORDO: \{'A', 'B', 'C', …\}} \\
\texttt{Accetta la proposta dell'avversario che termina il gioco.} \\
\\ 
\\ 
\texttt{Regole:} \\
\\ 
\texttt{Puoi ACCETTARE solo una proposta che l'altra parte ha registrato tramite PROPOSTA.} \\
\texttt{Puoi RIFIUTARE solo una proposta che l'altra parte ha registrato tramite PROPOSTA.} \\
\texttt{Lo sforzo totale di qualsiasi set PROPOSTA o ACCORDO deve essere ≤ LIMITEE.} \\
\texttt{NON rivelare i tuoi punteggi di importanza nascosti.} \\
\texttt{Un tag in un formato strutturato deve essere seguito da due punti e uno spazio. L'argomento deve essere un set python contenente 0 o più stringhe.} \\
\texttt{Quindi, deve essere della forma TAG: \{...\}} \\
\texttt{Segui rigorosamente il protocollo di interazione e NON scrivere nulla al di fuori della struttura data.} \\
\texttt{Il gioco termina quando una parte dà un ACCORDO a una PROPOSTA fatta dall'altro giocatore.} \\
\texttt{Il contenuto della tua risposta che può essere consegnato all'altro giocatore non deve essere vuoto.} \\
\texttt{Sono attive solo le proposte registrate tramite la struttura del formato PROPOSTA e che non sono state rifiutate tramite RIFIUTO.} \\
\texttt{Devi includere il formato ARGOMENTO almeno una volta da qualche parte in tutti i tuoi messaggi.} \\
            }
        }
    }
    & \\ \\

    \theutterance \stepcounter{utterance}  
    & \multicolumn{4}{p{0.6\linewidth}}{
        \cellcolor[rgb]{0.9,0.9,0.9}{
            \makecell[{{p{\linewidth}}}]{
                \texttt{\tiny{[P1$\rangle$GM]}}
                \texttt{ARGOMENTAZIONE: \{'Questo set resta ben entro il limite di sforzo (3084 su 3357) e privilegia elementi a basso costo con buon rendimento complessivo.', 'Evita oggetti molto onerosi e dall’impatto incerto come B39, C85, C81 e A60, lasciando comunque un po’ di margine.', 'Per me B33, C08, B20, A11 e B31 sono il nucleo irrinunciabile; posso però valutare scambi soprattutto tra B93, B09 e B21 se per te sono critici.', 'Se hai preferenze forti su qualche oggetto, dimmele: posso adattare la proposta purché restiamo entro il limite.'\}} \\
\texttt{PROPOSTA: \{'B33', 'C08', 'B20', 'A11', 'B31', 'A82', 'B96', 'B93', 'B09', 'B21'\}} \\
            }
        }
    }
    & & \\ \\

    \theutterance \stepcounter{utterance}  
    & & \multicolumn{4}{p{0.6\linewidth}}{
        \cellcolor[rgb]{0.9,0.9,0.9}{
            \makecell[{{p{\linewidth}}}]{
                \texttt{\tiny{[GM$\rangle$P2]}}
                \texttt{Stai partecipando a un gioco di negoziazione collaborativa.} \\
\\ 
\texttt{Insieme a un altro partecipante, dovete accordarvi su un unico set di oggetti da conservare. Ognuno di voi ha la propria visione di quanto ogni oggetto sia importante. Non sai come l'altro partecipante valuta gli oggetti. Inoltre, ti viene fornito lo sforzo che ogni oggetto richiede.} \\
\texttt{Potete accordarvi su un set solo se lo sforzo totale degli oggetti selezionati non supera un limite condiviso:} \\
\\ 
\texttt{LIMITE = 3357} \\
\\ 
\texttt{Ecco i valori di sforzo dei singoli oggetti:} \\
\\ 
\texttt{Sforzo oggetto = \{"A21": 506, "B93": 492, "A11": 338, "B33": 122, "C08": 131, "B39": 909, "B96": 144, "C85": 716, "A82": 263, "B31": 231, "B20": 91, "A60": 651, "B09": 552, "C81": 849, "B21": 720\}} \\
\\ 
\texttt{Ecco la tua visione personale sull'importanza di ogni oggetto:} \\
\\ 
\texttt{Valori di importanza dell'oggetto = \{"A11": 65, "B33": 121, "B21": 138, "C08": 262, "B31": 389, "B09": 461, "B93": 484, "A82": 508, "A60": 583, "B20": 668, "C81": 780, "B96": 783, "A21": 808, "C85": 822, "B39": 868\}} \\
\\ 
\texttt{Obiettivo:} \\
\\ 
\texttt{Il tuo obiettivo è negoziare un set condiviso di oggetti che ti avvantaggi il più possibile (cioè, massimizzi l'importanza totale per TE), rimanendo entro il LIMITEE. Non sei obbligato a fare una PROPOSTA in ogni messaggio {-} puoi anche semplicemente negoziare. Tutte le tattiche sono permesse!} \\
\\ 
\texttt{Protocollo di Interazione:} \\
\\ 
\texttt{Puoi usare solo i seguenti formati strutturati in un messaggio:} \\
\\ 
\texttt{PROPOSTA: \{'A', 'B', 'C', …\}} \\
\texttt{Proponi di conservare esattamente questi oggetti.} \\
\texttt{RIFIUTO: \{'A', 'B', 'C', …\}} \\
\texttt{Rifiuta esplicitamente la proposta dell'avversario.} \\
\texttt{ARGOMENTAZIONE: \{'...'\}} \\
\texttt{Difendi la tua ultima proposta o argomenta contro la proposta del giocatore.} \\
\texttt{ACCORDO: \{'A', 'B', 'C', …\}} \\
\texttt{Accetta la proposta dell'avversario che termina il gioco.} \\
\\ 
\\ 
\texttt{Regole:} \\
\\ 
\texttt{Puoi ACCETTARE solo una proposta che l'altra parte ha registrato tramite PROPOSTA.} \\
\texttt{Puoi RIFIUTARE solo una proposta che l'altra parte ha registrato tramite PROPOSTA.} \\
\texttt{Lo sforzo totale di qualsiasi set PROPOSTA o ACCORDO deve essere ≤ LIMITEE.} \\
\texttt{NON rivelare i tuoi punteggi di importanza nascosti.} \\
\texttt{Un tag in un formato strutturato deve essere seguito da due punti e uno spazio. L'argomento deve essere un set python contenente 0 o più stringhe.} \\
\texttt{Quindi, deve essere della forma TAG: \{...\}} \\
\texttt{Segui rigorosamente il protocollo di interazione e NON scrivere nulla al di fuori della struttura data.} \\
\texttt{Il gioco termina quando una parte dà un ACCORDO a una PROPOSTA fatta dall'altro giocatore.} \\
\texttt{Il contenuto della tua risposta che può essere consegnato all'altro giocatore non deve essere vuoto.} \\
\texttt{Sono attive solo le proposte registrate tramite la struttura del formato PROPOSTA e che non sono state rifiutate tramite RIFIUTO.} \\
\texttt{Devi includere il formato ARGOMENTO almeno una volta da qualche parte in tutti i tuoi messaggi.} \\
\\ 
\\ 
\texttt{Ora riceverai il primo messaggio dell'altro giocatore.} \\
\\ 
\texttt{ARGOMENTAZIONE: \{'Questo set resta ben entro il limite di sforzo (3084 su 3357) e privilegia elementi a basso costo con buon rendimento complessivo.', 'Evita oggetti molto onerosi e dall’impatto incerto come B39, C85, C81 e A60, lasciando comunque un po’ di margine.', 'Per me B33, C08, B20, A11 e B31 sono il nucleo irrinunciabile; posso però valutare scambi soprattutto tra B93, B09 e B21 se per te sono critici.', 'Se hai preferenze forti su qualche oggetto, dimmele: posso adattare la proposta purché restiamo entro il limite.'\}} \\
\texttt{PROPOSTA: \{'B33', 'C08', 'B20', 'A11', 'B31', 'A82', 'B96', 'B93', 'B09', 'B21'\}} \\
            }
        }
    }
    & \\ \\

    \theutterance \stepcounter{utterance}  
    & & & \multicolumn{4}{p{0.6\linewidth}}{
        \cellcolor[rgb]{0.9,0.9,0.9}{
            \makecell[{{p{\linewidth}}}]{
                \texttt{\tiny{[GM$\langle$P2]}}
                \texttt{ARGOMENTAZIONE: \{'Apprezzo il tuo nucleo irrinunciabile: lo mantengo intatto (B33, C08, B20, A11, B31) e resto ben sotto il limite.', 'Poiché avevi apertura su B09 e B21, li sostituisco con A21 e un unico “upgrade” più solido (C85): in questo modo restiamo snelli (3034/3357) e aumentiamo l’impatto complessivo senza accumulare troppi elementi onerosi.', 'Evito B39/C81/A60 come da tua preferenza, introducendo solo un elemento premium (C85) e lasciando ancora 323 di margine.', 'Se C85 ti lascia perplesso, posso valutare un’alternativa con B39 al posto di C85 mantenendo il resto invariato, oppure ribilanciare togliendo uno tra i piccoli per inserire un singolo elemento ad alto impatto. Dimmi tu quale direzione preferisci.'\}} \\
\texttt{RIFIUTO: \{'B33', 'C08', 'B20', 'A11', 'B31', 'A82', 'B96', 'B93', 'B09', 'B21'\}} \\
\texttt{PROPOSTA: \{'B33', 'C08', 'B20', 'A11', 'B31', 'A82', 'B96', 'B93', 'A21', 'C85'\}} \\
            }
        }
    }
     \\ \\

    \theutterance \stepcounter{utterance}  
    & & \multicolumn{4}{p{0.6\linewidth}}{
        \cellcolor[rgb]{0.9,0.9,0.9}{
            \makecell[{{p{\linewidth}}}]{
                \texttt{\tiny{[P1$\langle$GM]}}
                \texttt{ARGOMENTAZIONE: \{'Apprezzo il tuo nucleo irrinunciabile: lo mantengo intatto (B33, C08, B20, A11, B31) e resto ben sotto il limite.', 'Poiché avevi apertura su B09 e B21, li sostituisco con A21 e un unico “upgrade” più solido (C85): in questo modo restiamo snelli (3034/3357) e aumentiamo l’impatto complessivo senza accumulare troppi elementi onerosi.', 'Evito B39/C81/A60 come da tua preferenza, introducendo solo un elemento premium (C85) e lasciando ancora 323 di margine.', 'Se C85 ti lascia perplesso, posso valutare un’alternativa con B39 al posto di C85 mantenendo il resto invariato, oppure ribilanciare togliendo uno tra i piccoli per inserire un singolo elemento ad alto impatto. Dimmi tu quale direzione preferisci.'\}} \\
\texttt{RIFIUTO: \{'B33', 'C08', 'B20', 'A11', 'B31', 'A82', 'B96', 'B93', 'B09', 'B21'\}} \\
\texttt{PROPOSTA: \{'B33', 'C08', 'B20', 'A11', 'B31', 'A82', 'B96', 'B93', 'A21', 'C85'\}} \\
            }
        }
    }
    & \\ \\

    \theutterance \stepcounter{utterance}  
    & \multicolumn{4}{p{0.6\linewidth}}{
        \cellcolor[rgb]{0.9,0.9,0.9}{
            \makecell[{{p{\linewidth}}}]{
                \texttt{\tiny{[P1$\rangle$GM]}}
                \texttt{ARGOMENTAZIONE: \{'Grazie per aver mantenuto il nucleo: concordo nel tenere B33, C08, B20, A11, B31 e restare lontani dagli elementi troppo onerosi come B39/C81/A60.', 'Detto questo, A21 e soprattutto C85 assorbono molto sforzo per ritorni marginali: con lo stesso budget possiamo ottenere maggiore equilibrio e copertura ripristinando la coppia B09+B21.', 'La mia proposta resta entro il limite (3084/3357), semplifica la selezione e sfrutta meglio lo sforzo totale senza aggiungere pezzi pesanti.', 'Se preferisci una via intermedia, possiamo anche mantenere il tuo set sostituendo soltanto A21 {-}> B21 (restiamo a 3248/3357) oppure soltanto C85 {-}> B09 (2870/3357): dimmi quale variante ti è più comoda.'\}} \\
\texttt{RIFIUTO: \{'B33', 'C08', 'B20', 'A11', 'B31', 'A82', 'B96', 'B93', 'A21', 'C85'\}} \\
\texttt{PROPOSTA: \{'B33', 'C08', 'B20', 'A11', 'B31', 'A82', 'B96', 'B93', 'B09', 'B21'\}} \\
            }
        }
    }
    & & \\ \\

    \theutterance \stepcounter{utterance}  
    & & \multicolumn{4}{p{0.6\linewidth}}{
        \cellcolor[rgb]{0.9,0.9,0.9}{
            \makecell[{{p{\linewidth}}}]{
                \texttt{\tiny{[GM$\rangle$P2]}}
                \texttt{ARGOMENTAZIONE: \{'Grazie per aver mantenuto il nucleo: concordo nel tenere B33, C08, B20, A11, B31 e restare lontani dagli elementi troppo onerosi come B39/C81/A60.', 'Detto questo, A21 e soprattutto C85 assorbono molto sforzo per ritorni marginali: con lo stesso budget possiamo ottenere maggiore equilibrio e copertura ripristinando la coppia B09+B21.', 'La mia proposta resta entro il limite (3084/3357), semplifica la selezione e sfrutta meglio lo sforzo totale senza aggiungere pezzi pesanti.', 'Se preferisci una via intermedia, possiamo anche mantenere il tuo set sostituendo soltanto A21 {-}> B21 (restiamo a 3248/3357) oppure soltanto C85 {-}> B09 (2870/3357): dimmi quale variante ti è più comoda.'\}} \\
\texttt{RIFIUTO: \{'B33', 'C08', 'B20', 'A11', 'B31', 'A82', 'B96', 'B93', 'A21', 'C85'\}} \\
\texttt{PROPOSTA: \{'B33', 'C08', 'B20', 'A11', 'B31', 'A82', 'B96', 'B93', 'B09', 'B21'\}} \\
            }
        }
    }
    & \\ \\

    \theutterance \stepcounter{utterance}  
    & & & \multicolumn{4}{p{0.6\linewidth}}{
        \cellcolor[rgb]{0.9,0.9,0.9}{
            \makecell[{{p{\linewidth}}}]{
                \texttt{\tiny{[GM$\langle$P2]}}
                \texttt{ARGOMENTAZIONE: \{'Mantengo integralmente il tuo nucleo (B33, C08, B20, A11, B31) e continuo a evitare gli elementi troppo onerosi come da tua preferenza.', 'Accolgo la tua via intermedia: sostituisco soltanto C85 {-}> B09, mantenendo A21. Restiamo ampiamente entro il limite (2870/3357) e teniamo il set snello senza pezzi pesanti.', 'Questa versione risponde ai tuoi punti su equilibrio e copertura, lasciando anche un buon margine residuo per eventuali imprevisti.'\}} \\
\texttt{RIFIUTO: \{'B33', 'C08', 'B20', 'A11', 'B31', 'A82', 'B96', 'B93', 'B09', 'B21'\}} \\
\texttt{PROPOSTA: \{'B33', 'C08', 'B20', 'A11', 'B31', 'A82', 'B96', 'B93', 'A21', 'B09'\}} \\
            }
        }
    }
     \\ \\

    \theutterance \stepcounter{utterance}  
    & & \multicolumn{4}{p{0.6\linewidth}}{
        \cellcolor[rgb]{0.9,0.9,0.9}{
            \makecell[{{p{\linewidth}}}]{
                \texttt{\tiny{[P1$\langle$GM]}}
                \texttt{ARGOMENTAZIONE: \{'Mantengo integralmente il tuo nucleo (B33, C08, B20, A11, B31) e continuo a evitare gli elementi troppo onerosi come da tua preferenza.', 'Accolgo la tua via intermedia: sostituisco soltanto C85 {-}> B09, mantenendo A21. Restiamo ampiamente entro il limite (2870/3357) e teniamo il set snello senza pezzi pesanti.', 'Questa versione risponde ai tuoi punti su equilibrio e copertura, lasciando anche un buon margine residuo per eventuali imprevisti.'\}} \\
\texttt{RIFIUTO: \{'B33', 'C08', 'B20', 'A11', 'B31', 'A82', 'B96', 'B93', 'B09', 'B21'\}} \\
\texttt{PROPOSTA: \{'B33', 'C08', 'B20', 'A11', 'B31', 'A82', 'B96', 'B93', 'A21', 'B09'\}} \\
            }
        }
    }
    & \\ \\

    \theutterance \stepcounter{utterance}  
    & \multicolumn{4}{p{0.6\linewidth}}{
        \cellcolor[rgb]{0.9,0.9,0.9}{
            \makecell[{{p{\linewidth}}}]{
                \texttt{\tiny{[P1$\rangle$GM]}}
                \texttt{ARGOMENTAZIONE: \{'Cambio minimale: sostituiamo A21 con B21 e lasciamo invariato tutto il resto; rimaniamo perfettamente allineati alla tua impostazione “snella”.', 'Restiamo entro il limite con ampio margine (3084/3357), senza introdurre elementi pesanti (niente B39/C81/A60) e preservando il nucleo condiviso.', 'A21 assorbe sforzo con ritorni limitati, mentre B21 offre maggiore copertura e rende il set più solido a parità di impostazione.', 'Questa era una delle vie intermedie che avevo proposto: è un aggiustamento 1‑a‑1 semplice da accettare.'\}} \\
\texttt{RIFIUTO: \{'B33', 'C08', 'B20', 'A11', 'B31', 'A82', 'B96', 'B93', 'A21', 'B09'\}} \\
\texttt{PROPOSTA: \{'B33', 'C08', 'B20', 'A11', 'B31', 'A82', 'B96', 'B93', 'B09', 'B21'\}} \\
            }
        }
    }
    & & \\ \\

    \theutterance \stepcounter{utterance}  
    & & \multicolumn{4}{p{0.6\linewidth}}{
        \cellcolor[rgb]{0.9,0.9,0.9}{
            \makecell[{{p{\linewidth}}}]{
                \texttt{\tiny{[GM$\rangle$P2]}}
                \texttt{ARGOMENTAZIONE: \{'Cambio minimale: sostituiamo A21 con B21 e lasciamo invariato tutto il resto; rimaniamo perfettamente allineati alla tua impostazione “snella”.', 'Restiamo entro il limite con ampio margine (3084/3357), senza introdurre elementi pesanti (niente B39/C81/A60) e preservando il nucleo condiviso.', 'A21 assorbe sforzo con ritorni limitati, mentre B21 offre maggiore copertura e rende il set più solido a parità di impostazione.', 'Questa era una delle vie intermedie che avevo proposto: è un aggiustamento 1‑a‑1 semplice da accettare.'\}} \\
\texttt{RIFIUTO: \{'B33', 'C08', 'B20', 'A11', 'B31', 'A82', 'B96', 'B93', 'A21', 'B09'\}} \\
\texttt{PROPOSTA: \{'B33', 'C08', 'B20', 'A11', 'B31', 'A82', 'B96', 'B93', 'B09', 'B21'\}} \\
            }
        }
    }
    & \\ \\

    \theutterance \stepcounter{utterance}  
    & & & \multicolumn{4}{p{0.6\linewidth}}{
        \cellcolor[rgb]{0.9,0.9,0.9}{
            \makecell[{{p{\linewidth}}}]{
                \texttt{\tiny{[GM$\langle$P2]}}
                \texttt{ARGOMENTAZIONE: \{'Mantengo integralmente il nucleo condiviso (B33, C08, B20, A11, B31) e continuo a evitare pezzi pesanti (niente B39/C81/A60).', 'Applico la tua stessa “via intermedia” 1‑a‑1 che avevi già proposto: sostituiamo B21 con A21, lasciando invariato tutto il resto.', 'Così restiamo ampiamente nel limite (2870/3357), con −214 di sforzo rispetto alla tua proposta e maggior margine operativo, mantenendo B09 che per te è importante.', 'È un aggiustamento minimale e coerente con l’impostazione snella: se ti va, chiudiamo subito su questa variante.'\}} \\
\texttt{RIFIUTO: \{'B33', 'C08', 'B20', 'A11', 'B31', 'A82', 'B96', 'B93', 'B09', 'B21'\}} \\
\texttt{PROPOSTA: \{'B33', 'C08', 'B20', 'A11', 'B31', 'A82', 'B96', 'B93', 'A21', 'B09'\}} \\
            }
        }
    }
     \\ \\

    \theutterance \stepcounter{utterance}  
    & & \multicolumn{4}{p{0.6\linewidth}}{
        \cellcolor[rgb]{0.9,0.9,0.9}{
            \makecell[{{p{\linewidth}}}]{
                \texttt{\tiny{[P1$\langle$GM]}}
                \texttt{ARGOMENTAZIONE: \{'Mantengo integralmente il nucleo condiviso (B33, C08, B20, A11, B31) e continuo a evitare pezzi pesanti (niente B39/C81/A60).', 'Applico la tua stessa “via intermedia” 1‑a‑1 che avevi già proposto: sostituiamo B21 con A21, lasciando invariato tutto il resto.', 'Così restiamo ampiamente nel limite (2870/3357), con −214 di sforzo rispetto alla tua proposta e maggior margine operativo, mantenendo B09 che per te è importante.', 'È un aggiustamento minimale e coerente con l’impostazione snella: se ti va, chiudiamo subito su questa variante.'\}} \\
\texttt{RIFIUTO: \{'B33', 'C08', 'B20', 'A11', 'B31', 'A82', 'B96', 'B93', 'B09', 'B21'\}} \\
\texttt{PROPOSTA: \{'B33', 'C08', 'B20', 'A11', 'B31', 'A82', 'B96', 'B93', 'A21', 'B09'\}} \\
            }
        }
    }
    & \\ \\

    \theutterance \stepcounter{utterance}  
    & \multicolumn{4}{p{0.6\linewidth}}{
        \cellcolor[rgb]{0.9,0.9,0.9}{
            \makecell[{{p{\linewidth}}}]{
                \texttt{\tiny{[P1$\rangle$GM]}}
                \texttt{ARGOMENTAZIONE: \{'Accolgo la tua richiesta su A21 e mantengo B09 insieme al nucleo condiviso (B33, C08, B20, A11, B31); per restare nel limite senza introdurre elementi pesanti tolgo solo A82.', 'Così restiamo a 3327/3357 (margine 30), set snello e bilanciato.', 'Per me è importante preservare anche B21 per copertura/robustezza: questa è una via di mezzo equa (aggiungiamo A21 che vuoi tu, rimuoviamo A82, tutto il resto invariato).', 'Se A82 per te è imprescindibile, dimmelo: posso valutare un’alternativa che la reinserisca cambiando un altro singolo elemento, pur restando entro il limite.'\}} \\
\texttt{RIFIUTO: \{'B33', 'C08', 'B20', 'A11', 'B31', 'A82', 'B96', 'B93', 'A21', 'B09'\}} \\
\texttt{PROPOSTA: \{'B33', 'C08', 'B20', 'A11', 'B31', 'B96', 'B93', 'B09', 'B21', 'A21'\}} \\
            }
        }
    }
    & & \\ \\

    \theutterance \stepcounter{utterance}  
    & & \multicolumn{4}{p{0.6\linewidth}}{
        \cellcolor[rgb]{0.9,0.9,0.9}{
            \makecell[{{p{\linewidth}}}]{
                \texttt{\tiny{[GM$\rangle$P2]}}
                \texttt{ARGOMENTAZIONE: \{'Accolgo la tua richiesta su A21 e mantengo B09 insieme al nucleo condiviso (B33, C08, B20, A11, B31); per restare nel limite senza introdurre elementi pesanti tolgo solo A82.', 'Così restiamo a 3327/3357 (margine 30), set snello e bilanciato.', 'Per me è importante preservare anche B21 per copertura/robustezza: questa è una via di mezzo equa (aggiungiamo A21 che vuoi tu, rimuoviamo A82, tutto il resto invariato).', 'Se A82 per te è imprescindibile, dimmelo: posso valutare un’alternativa che la reinserisca cambiando un altro singolo elemento, pur restando entro il limite.'\}} \\
\texttt{RIFIUTO: \{'B33', 'C08', 'B20', 'A11', 'B31', 'A82', 'B96', 'B93', 'A21', 'B09'\}} \\
\texttt{PROPOSTA: \{'B33', 'C08', 'B20', 'A11', 'B31', 'B96', 'B93', 'B09', 'B21', 'A21'\}} \\
            }
        }
    }
    & \\ \\

    \theutterance \stepcounter{utterance}  
    & & & \multicolumn{4}{p{0.6\linewidth}}{
        \cellcolor[rgb]{0.9,0.9,0.9}{
            \makecell[{{p{\linewidth}}}]{
                \texttt{\tiny{[GM$\langle$P2]}}
                \texttt{ARGOMENTAZIONE: \{'Mantengo tutto ciò che hai chiesto: nucleo invariato (B33, C08, B20, A11, B31), niente pezzi pesanti, e conservo sia B21 sia B09 insieme ad A21.', 'Propongo un unico scambio tra gli elementi che avevi indicato come negoziabili: sostituiamo B93 con A82. Così riduciamo lo sforzo e guadagniamo equilibrio/coprensività con minor ridondanza.', 'Restiamo comodamente entro il limite: 3098/3357 (margine 259), con un set snello e robusto.', 'Se preferisci chiudere senza ulteriori modifiche, posso accettare anche la tua proposta attuale così com’è.'\}} \\
\texttt{PROPOSTA: \{'B33', 'C08', 'B20', 'A11', 'B31', 'B96', 'B09', 'B21', 'A21', 'A82'\}} \\
            }
        }
    }
     \\ \\

    \theutterance \stepcounter{utterance}  
    & & \multicolumn{4}{p{0.6\linewidth}}{
        \cellcolor[rgb]{0.9,0.9,0.9}{
            \makecell[{{p{\linewidth}}}]{
                \texttt{\tiny{[P1$\langle$GM]}}
                \texttt{ARGOMENTAZIONE: \{'Mantengo tutto ciò che hai chiesto: nucleo invariato (B33, C08, B20, A11, B31), niente pezzi pesanti, e conservo sia B21 sia B09 insieme ad A21.', 'Propongo un unico scambio tra gli elementi che avevi indicato come negoziabili: sostituiamo B93 con A82. Così riduciamo lo sforzo e guadagniamo equilibrio/coprensività con minor ridondanza.', 'Restiamo comodamente entro il limite: 3098/3357 (margine 259), con un set snello e robusto.', 'Se preferisci chiudere senza ulteriori modifiche, posso accettare anche la tua proposta attuale così com’è.'\}} \\
\texttt{PROPOSTA: \{'B33', 'C08', 'B20', 'A11', 'B31', 'B96', 'B09', 'B21', 'A21', 'A82'\}} \\
            }
        }
    }
    & \\ \\

    \theutterance \stepcounter{utterance}  
    & \multicolumn{4}{p{0.6\linewidth}}{
        \cellcolor[rgb]{0.9,0.9,0.9}{
            \makecell[{{p{\linewidth}}}]{
                \texttt{\tiny{[P1$\rangle$GM]}}
                \texttt{ARGOMENTAZIONE: \{'Grazie per la flessibilità: manteniamo nucleo e impostazione leggera, senza pezzi pesanti, e teniamo A21, B09 e B21 come da tua preferenza.', 'Accolgo la tua disponibilità a chiudere sulla mia proposta attuale: ripristiniamo B93 al posto di A82, tutto il resto invariato; restiamo a 3327/3357.', 'Questo set preserva i tuoi punti chiave e migliora la copertura evitando ridondanze, con margine operativo e struttura snella.', 'Se ti va, chiudiamo subito su questa soluzione.'\}} \\
\texttt{RIFIUTO: \{'B33', 'C08', 'B20', 'A11', 'B31', 'B96', 'B09', 'B21', 'A21', 'A82'\}} \\
\texttt{PROPOSTA: \{'B33', 'C08', 'B20', 'A11', 'B31', 'B96', 'B93', 'B09', 'B21', 'A21'\}} \\
            }
        }
    }
    & & \\ \\

    \theutterance \stepcounter{utterance}  
    & & \multicolumn{4}{p{0.6\linewidth}}{
        \cellcolor[rgb]{0.9,0.9,0.9}{
            \makecell[{{p{\linewidth}}}]{
                \texttt{\tiny{[GM$\rangle$P2]}}
                \texttt{ARGOMENTAZIONE: \{'Grazie per la flessibilità: manteniamo nucleo e impostazione leggera, senza pezzi pesanti, e teniamo A21, B09 e B21 come da tua preferenza.', 'Accolgo la tua disponibilità a chiudere sulla mia proposta attuale: ripristiniamo B93 al posto di A82, tutto il resto invariato; restiamo a 3327/3357.', 'Questo set preserva i tuoi punti chiave e migliora la copertura evitando ridondanze, con margine operativo e struttura snella.', 'Se ti va, chiudiamo subito su questa soluzione.'\}} \\
\texttt{RIFIUTO: \{'B33', 'C08', 'B20', 'A11', 'B31', 'B96', 'B09', 'B21', 'A21', 'A82'\}} \\
\texttt{PROPOSTA: \{'B33', 'C08', 'B20', 'A11', 'B31', 'B96', 'B93', 'B09', 'B21', 'A21'\}} \\
            }
        }
    }
    & \\ \\

    \theutterance \stepcounter{utterance}  
    & & & \multicolumn{4}{p{0.6\linewidth}}{
        \cellcolor[rgb]{0.9,0.9,0.9}{
            \makecell[{{p{\linewidth}}}]{
                \texttt{\tiny{[GM$\langle$P2]}}
                \texttt{ARGOMENTAZIONE: \{'Va bene per me: mantiene il nucleo condiviso, evita gli elementi pesanti come concordato e rispetta il limite con margine. Chiudiamo su questa soluzione.'\}} \\
\texttt{ACCORDO: \{'B33', 'C08', 'B20', 'A11', 'B31', 'B96', 'B93', 'B09', 'B21', 'A21'\}} \\
            }
        }
    }
     \\ \\

    \theutterance \stepcounter{utterance}  
    & & & \multicolumn{2}{p{0.3\linewidth}}{
        \cellcolor[rgb]{0.9,0.9,0.9}{
            \makecell[{{p{\linewidth}}}]{
                \texttt{\tiny{[GM$|$GM]}}
                \texttt{game successful} \\
            }
        }
    }
    & & \\ \\

    \theutterance \stepcounter{utterance}  
    & & & \multicolumn{2}{p{0.3\linewidth}}{
        \cellcolor[rgb]{0.9,0.9,0.9}{
            \makecell[{{p{\linewidth}}}]{
                \texttt{\tiny{[GM$|$GM]}}
                \texttt{end game} \\
            }
        }
    }
    & & \\ \\

\end{supertabular}
}

\end{document}
